% Options for packages loaded elsewhere
\PassOptionsToPackage{unicode}{hyperref}
\PassOptionsToPackage{hyphens}{url}
\documentclass[
]{article}
\usepackage{xcolor}
\usepackage[margin=1in]{geometry}
\usepackage{amsmath,amssymb}
\setcounter{secnumdepth}{-\maxdimen} % remove section numbering
\usepackage{iftex}
\ifPDFTeX
  \usepackage[T1]{fontenc}
  \usepackage[utf8]{inputenc}
  \usepackage{textcomp} % provide euro and other symbols
\else % if luatex or xetex
  \usepackage{unicode-math} % this also loads fontspec
  \defaultfontfeatures{Scale=MatchLowercase}
  \defaultfontfeatures[\rmfamily]{Ligatures=TeX,Scale=1}
\fi
\usepackage{lmodern}
\ifPDFTeX\else
  % xetex/luatex font selection
\fi
% Use upquote if available, for straight quotes in verbatim environments
\IfFileExists{upquote.sty}{\usepackage{upquote}}{}
\IfFileExists{microtype.sty}{% use microtype if available
  \usepackage[]{microtype}
  \UseMicrotypeSet[protrusion]{basicmath} % disable protrusion for tt fonts
}{}
\makeatletter
\@ifundefined{KOMAClassName}{% if non-KOMA class
  \IfFileExists{parskip.sty}{%
    \usepackage{parskip}
  }{% else
    \setlength{\parindent}{0pt}
    \setlength{\parskip}{6pt plus 2pt minus 1pt}}
}{% if KOMA class
  \KOMAoptions{parskip=half}}
\makeatother
\usepackage{color}
\usepackage{fancyvrb}
\newcommand{\VerbBar}{|}
\newcommand{\VERB}{\Verb[commandchars=\\\{\}]}
\DefineVerbatimEnvironment{Highlighting}{Verbatim}{commandchars=\\\{\}}
% Add ',fontsize=\small' for more characters per line
\usepackage{framed}
\definecolor{shadecolor}{RGB}{248,248,248}
\newenvironment{Shaded}{\begin{snugshade}}{\end{snugshade}}
\newcommand{\AlertTok}[1]{\textcolor[rgb]{0.94,0.16,0.16}{#1}}
\newcommand{\AnnotationTok}[1]{\textcolor[rgb]{0.56,0.35,0.01}{\textbf{\textit{#1}}}}
\newcommand{\AttributeTok}[1]{\textcolor[rgb]{0.13,0.29,0.53}{#1}}
\newcommand{\BaseNTok}[1]{\textcolor[rgb]{0.00,0.00,0.81}{#1}}
\newcommand{\BuiltInTok}[1]{#1}
\newcommand{\CharTok}[1]{\textcolor[rgb]{0.31,0.60,0.02}{#1}}
\newcommand{\CommentTok}[1]{\textcolor[rgb]{0.56,0.35,0.01}{\textit{#1}}}
\newcommand{\CommentVarTok}[1]{\textcolor[rgb]{0.56,0.35,0.01}{\textbf{\textit{#1}}}}
\newcommand{\ConstantTok}[1]{\textcolor[rgb]{0.56,0.35,0.01}{#1}}
\newcommand{\ControlFlowTok}[1]{\textcolor[rgb]{0.13,0.29,0.53}{\textbf{#1}}}
\newcommand{\DataTypeTok}[1]{\textcolor[rgb]{0.13,0.29,0.53}{#1}}
\newcommand{\DecValTok}[1]{\textcolor[rgb]{0.00,0.00,0.81}{#1}}
\newcommand{\DocumentationTok}[1]{\textcolor[rgb]{0.56,0.35,0.01}{\textbf{\textit{#1}}}}
\newcommand{\ErrorTok}[1]{\textcolor[rgb]{0.64,0.00,0.00}{\textbf{#1}}}
\newcommand{\ExtensionTok}[1]{#1}
\newcommand{\FloatTok}[1]{\textcolor[rgb]{0.00,0.00,0.81}{#1}}
\newcommand{\FunctionTok}[1]{\textcolor[rgb]{0.13,0.29,0.53}{\textbf{#1}}}
\newcommand{\ImportTok}[1]{#1}
\newcommand{\InformationTok}[1]{\textcolor[rgb]{0.56,0.35,0.01}{\textbf{\textit{#1}}}}
\newcommand{\KeywordTok}[1]{\textcolor[rgb]{0.13,0.29,0.53}{\textbf{#1}}}
\newcommand{\NormalTok}[1]{#1}
\newcommand{\OperatorTok}[1]{\textcolor[rgb]{0.81,0.36,0.00}{\textbf{#1}}}
\newcommand{\OtherTok}[1]{\textcolor[rgb]{0.56,0.35,0.01}{#1}}
\newcommand{\PreprocessorTok}[1]{\textcolor[rgb]{0.56,0.35,0.01}{\textit{#1}}}
\newcommand{\RegionMarkerTok}[1]{#1}
\newcommand{\SpecialCharTok}[1]{\textcolor[rgb]{0.81,0.36,0.00}{\textbf{#1}}}
\newcommand{\SpecialStringTok}[1]{\textcolor[rgb]{0.31,0.60,0.02}{#1}}
\newcommand{\StringTok}[1]{\textcolor[rgb]{0.31,0.60,0.02}{#1}}
\newcommand{\VariableTok}[1]{\textcolor[rgb]{0.00,0.00,0.00}{#1}}
\newcommand{\VerbatimStringTok}[1]{\textcolor[rgb]{0.31,0.60,0.02}{#1}}
\newcommand{\WarningTok}[1]{\textcolor[rgb]{0.56,0.35,0.01}{\textbf{\textit{#1}}}}
\usepackage{graphicx}
\makeatletter
\newsavebox\pandoc@box
\newcommand*\pandocbounded[1]{% scales image to fit in text height/width
  \sbox\pandoc@box{#1}%
  \Gscale@div\@tempa{\textheight}{\dimexpr\ht\pandoc@box+\dp\pandoc@box\relax}%
  \Gscale@div\@tempb{\linewidth}{\wd\pandoc@box}%
  \ifdim\@tempb\p@<\@tempa\p@\let\@tempa\@tempb\fi% select the smaller of both
  \ifdim\@tempa\p@<\p@\scalebox{\@tempa}{\usebox\pandoc@box}%
  \else\usebox{\pandoc@box}%
  \fi%
}
% Set default figure placement to htbp
\def\fps@figure{htbp}
\makeatother
\setlength{\emergencystretch}{3em} % prevent overfull lines
\providecommand{\tightlist}{%
  \setlength{\itemsep}{0pt}\setlength{\parskip}{0pt}}
\usepackage{bookmark}
\IfFileExists{xurl.sty}{\usepackage{xurl}}{} % add URL line breaks if available
\urlstyle{same}
\hypersetup{
  pdftitle={Coding Challenge 4},
  pdfauthor={Kamerin Sandlin},
  hidelinks,
  pdfcreator={LaTeX via pandoc}}

\title{Coding Challenge 4}
\author{Kamerin Sandlin}
\date{2026-02-19}

\begin{document}
\maketitle

{
\setcounter{tocdepth}{2}
\tableofcontents
}
Coding Challenge 4

\section{Question 1a}\label{question-1a}

AA YAML Header is the title information at the top of a document. It can
display an HTML file by default which cannot be read on GitHub.

\section{Question 1b}\label{question-1b}

Literate programming is pretty much exactly what I'm doing. Making
source code in a repository but annotating it in plain English so that
it is explained what the code does.

\section{Question 2}\label{question-2}

Take the code you wrote for coding challenge 3, question 5, and
incorporate it into your R Markdown file.

\subsection{2a}\label{a}

At the top of the document, make a clickable link to the manuscript
where these data are published.

\subsection{2b}\label{b}

Read the data using a relative file path with \texttt{na.strings} option
set to ``na''.

\subsection{2c}\label{c}

Make a separate code chunk for the figures plotting the DON data,
15ADON, and Seedmass, and one for the three combined using
\texttt{ggarrange}.

\section{Question 3}\label{question-3}

Knit your document together into ONE format of choice (docx, pdf, md).

\section{Question 4}\label{question-4}

Push the .docx or .pdf and .md files to GitHub inside a directory called
Coding Challenge 4.

\section{Question 5}\label{question-5}

Now edit, commit, and push the README file for your repository and
include the following elements:

\#Question 1a AA YAML Header is the title information at the top of a
document. It can display an HTML file by default which cannot be read on
github.

\#Question 1B Literate programming is pretty much exactly what I'm
doing. Making source code in a repository but annotating it in plain
english to where it is explained what that code does.

\#Question 2 Take the code you wrote for coding challenge 3, question 5,
and incorporate it into your R markdown file.

ANSWER: SeedMassgwc \textless- SeedMassPlot + geom\_pwc(data = MycData,
aes(group = Treatment), method = ``t.test'', label =
``\{p.adj.format\}\{p.adj.signif\}'') This is how you add the pairwise
comparison to a plot, seed mass plot for example.

FinalPlot \textless- ggarrange(SeedMassgwc, X15gwc, DONgwcPlot, ncol =
3, nrow = 1, common.legend = TRUE) This is how the final plot was all
put together using ggarrange.

\#Question 2A At the top of the document, make a clickable link to the
manuscript where these data are published.

ANSWER: Noel, Z.A., Roze, L.V., Breunig, M., Trail, F. 2022. Endophytic
fungi as promising biocontrol agent to protect wheat from Fusarium
graminearum head blight. Plant Disease.
\url{https://doi.org/10.1094/PDIS-06-21-1253-RE}

\#Question 2B Read the data using a relative file path with na.strings
option set to ``na''. This means you need to put the Mycotoxin.csv file
we have used for the past two weeks into your directory, which git
tracks.

ANSWER: MycotoxinData\('15ADON'[MycotoxinData\)'15ADON` \%in\% c(``nd'',
``na''){]} \textless- NA This is because the 15ADON column was the only
one giving na.strings troubles, so the ``na.strings'' code was mapped to
15ADON specifically.

\#Question 2C Make a separate code chunk for the figures plotting the
DON data, 15ADON, and Seedmass, and one for the three combined using
ggarrange.

Library ggplot2

\begin{Shaded}
\begin{Highlighting}[]
\FunctionTok{library}\NormalTok{(ggplot2)}
\FunctionTok{library}\NormalTok{(ggpubr)}
\end{Highlighting}
\end{Shaded}

\begin{Shaded}
\begin{Highlighting}[]
\CommentTok{\#DON Plot}
\CommentTok{\#ggplot(MycData, aes(x = Treatment, y = DON, fill = Cultivar)) + geom\_boxplot(position = position\_dodge(width = 0.75)) + geom\_jitter(aes(color = Cultivar), position = position\_jitterdodge(jitter.width = 0.2, dodge.width = 1), shape = 20, size = 2, alpha = 0.6) + scale\_fill\_manual(values = mycolors) + scale\_color\_manual(values = mycolors) + ylab("DON (ppm)") + xlab("") + theme\_classic() + coord\_cartesian(ylim = c(0, 210)) + facet\_wrap(\textasciitilde{}Cultivar)}

\CommentTok{\#X15Plot}
\CommentTok{\#ggplot(MycotoxinData, aes(x = Treatment, y = \textquotesingle{}15ADON\textquotesingle{}, fill = Cultivar)) + geom\_boxplot(position = position\_dodge(width = 0.75)) + geom\_jitter(aes(color = Cultivar), position = position\_jitterdodge(jitter.width = 0.2, dodge.width = 0.9), alpha = 0.6, size = 2) + scale\_fill\_manual(values = mycolors) + scale\_color\_manual(values = mycolors) + ylab("15ADON") + xlab("") + theme\_classic()}

\CommentTok{\#Seed Mass Plot}
\CommentTok{\#ggplot(MycData, aes(x = Treatment, y = MassperSeed\_mg, fill = Cultivar)) + geom\_boxplot(position = position\_dodge(width = 0.9)) + geom\_jitter(aes(color = Cultivar), position = position\_jitterdodge(jitter.width = 0.2, dodge.width = 0.9), alpha = 0.6, size = 2) + scale\_fill\_manual(values = mycolors) + scale\_color\_manual(values = mycolors) + ylab("Seed Mass (mg)") + xlab("") + theme\_classic() }

\CommentTok{\#Combined Plot, using ggarrange}
\CommentTok{\#CombinedPlot \textless{}{-} ggarrange(DONPlot, SeedMassPlot, X15Plot, labels =c("A", "B", "C"), common.legend = TRUE) }
\end{Highlighting}
\end{Shaded}

\#Question 3 \#Knit your document together into ONE format of choice
(docx, pdf, md) Done via knit option, I picked HTML file type.

\#Question 4 Push the .docx or .pdf and .md files to GitHub inside a
directory called Coding Challenge 4. Pushed.

\#Question 5 Now edit, commit, and push the README file for your
repository and include the following elements: A. A clickable link in
your README to your GitHub flavored .md file B. A file tree of your
GitHub repository.

\end{document}
